\documentclass[conference]{IEEEtran}

\usepackage{amsmath,amssymb,amsfonts}
\usepackage{graphicx}
\usepackage{booktabs}
\usepackage{textcomp}
\usepackage{xcolor}
\usepackage{algorithm}
\usepackage{algpseudocode}
\usepackage{url}
\usepackage{cite}
\usepackage{multirow}
\usepackage{array}

\begin{document}

\title{MTD-HealthNet: A Zone-Based Moving Target Defense Framework with Dynamic NAT Randomization for Securing Hospital Local Area Networks}

\author{
\IEEEauthorblockN{Rushyendra Reddy Burla\textsuperscript{1}, Kishan Sai Yallanuru\textsuperscript{2},\\ Satya Siddardha Vajjula\textsuperscript{3}, Jeeshitha Karri\textsuperscript{4}}
\IEEEauthorblockA{
Department of Computer Science and Engineering (Cyber Security) \\
TIFAC-CORE in Cyber Security, Amrita School of Computing \\
Amrita Vishwa Vidyapeetham, Coimbatore, India \\
\textsuperscript{1}cb.en.u4cys22015@cb.students.amrita.edu,
\textsuperscript{2}cb.en.u4cys22064@cb.students.amrita.edu \\
\textsuperscript{3}cb.en.u4cys22063@cb.students.amrita.edu,
\textsuperscript{4}cb.en.u4cys22032@cb.students.amrita.edu
}
}

\maketitle

\begin{abstract}
Hospital devices usually sit at the same IP for months, which gives attackers plenty of time to map everything out. We built \textbf{MTD-HealthNet}, an SDN controller that keeps shuffling the external IPs of hospital machines while their internal addresses stay the same. Three zones---Critical (120\,s), Imaging (100\,s), Administrative (80\,s)---rotate at different speeds depending on clinical importance. Internally hosts use fixed 10.0.0.0/24 addresses; only the 172.16.0.0/16 facing address changes. On a six-host Mininet testbed, scans could only find \textbf{1 of 6} hosts on average once rotation was active, which is \textbf{83.3\%} worse than a static network for the attacker. Scanning effort went up by \textbf{4.7$\times$}. Legitimate traffic still delivered at \textbf{97.8\%}, and each address swap finished in about \textbf{12.4~ms}.
\end{abstract}

\begin{IEEEkeywords}
Moving Target Defense, Hospital Network Security, Software-Defined Networking, IP Randomization, NAT Translation, Healthcare Cybersecurity, OpenFlow, Cyber Resilience
\end{IEEEkeywords}

\section{Introduction}

Inside any hospital, hundreds of devices share one network. Monitors push vitals to nursing stations. Imaging workstations fetch scans from archives. Pharmacy apps check prescriptions against patient records. If one machine gets compromised, the rest are reachable over the same wires.

Why is that so dangerous? Because IPs basicaly never change. A machine gets an address during setup and keeps it for years. One network scan gives an attacker a map of everything, and that map stays accurate indefinitely~\cite{lei2018survey, cai2016mtd}. Ransomware, worms, insider threats---they all rely on this predictability.

Conventional tools like firewalls and IDS try to protect a network that stays still. They add walls around a fixed layout but dont move anything. An attacker who already mapped the place can take their time probing for holes in those walls.

A diferent approach is to keep changing the layout. When addresses shift every couple of minutes, old scan results expire fast. We designed \textbf{MTD-HealthNet} around this idea, with four goals that set it apart from earlier work:

\begin{itemize}
    \item \textbf{Zone-aware rotation.} Devices grouped by clinical risk rotate at different rates. Life-support gear gets a slow, stable cycle. Office computers rotate aggressively since a brief blip there has no clinical consequence.
    \item \textbf{Dual-address scheme.} Each device holds a permanent internal IP. Only the external-facing address rotates. No application changes needed.
    \item \textbf{Application-level validation.} We tested actual encrypted data transfers between hosts, not just ICMP pings.
    \item \textbf{Quantified trade-offs.} We measured both what the attacker loses and what the hospital pays in performace overhead.
\end{itemize}

\section{Related Work}

\subsection{MTD Landscape}

Moving Target Defense covers a range of mutation strategies at different layers~\cite{lei2018survey}. Network-layer techniques change IPs, ports, or routes. Host-layer ones randomize memory layouts or OS fingerprints. Application-layer approaches diversify configs or software stacks. Our work sits at the network layer since thats where scanning happens.

A well-known tension exists here: faster mutations degrade attacker capability but also put load on the defender's infrastucture~\cite{cai2016mtd}. Balancing that is harder when the network carries time-sensitive medical data.

\subsection{OpenFlow-Based Address Mutation}

Programmable switches can reassign addresses without hurting throughput, as demonstrated in~\cite{dishington2019ipmux}. Automated scans became less accurate. But their testbed gave every host the same treatment. Hospitals cant do that---a ventilator and a billing PC have very different tolerance for disruption.

A follow-up study~\cite{zal2024lossless} addressed the brief packet loss during rule transitions by keeping old and new forwarding entries active simultaneously for a short overlap. We adopted a similar dual-rule window.

\subsection{Legacy Device Compatibility}

Hospital gear often runs firmware that cant be patched. Huang et al.~\cite{huang2021legacy} handled this by performing address changes outside the device---it keeps its factory IP while a proxy layer does the translation. Our dual-address model is based on the same idea.

\subsection{Gaps in Existing Work}

We identified four issues that existing literature hadnt addressed together:

\begin{enumerate}
    \item Rotation speed is usually the same for all devices, ignoring that different categories tolerate disruption differently.
    \item Most evaluations only use ICMP. That cant confirm if actual application data survives a rotation.
    \item No prior system combined zone-based access policies with IP rotation specifically for healthcare.
    \item There are no benchmarks capturing both security improvements and operational costs for hospital-like setups.
\end{enumerate}

MTD-HealthNet was built to fill all four. Table~\ref{tab:litcomp} compares the main mutation approaches.

\begin{table}[htbp]
\centering
\caption{Comparison of Moving Target Defense Techniques}
\label{tab:litcomp}
\begin{tabular}{p{1.4cm}p{2cm}p{2cm}p{2cm}}
\toprule
\textbf{Technique} & \textbf{Security Effectiveness} & \textbf{Overhead} & \textbf{Deployment Issues} \\
\midrule
IP Hopping~\cite{dishington2019ipmux, zal2024lossless} & High against reconnaissance and worm propagation & Requires coordinated address management & May challenge session continuity for legacy applications \\
\midrule
Port Hopping~\cite{cai2016mtd} & High against port-based attacks & Requires client-server synchronization & May disrupt long-lived sessions; needs protocol adaptation \\
\midrule
NAT Randomization~\cite{cai2016mtd} & Provides additional protection at NAT boundaries & Moderate complexity; customized NAT management & Needs careful timer and state management for auditability \\
\bottomrule
\end{tabular}
\end{table}

\section{System Architecture}

The framework has three layers with clear boundaries. Each can be modified or replaced without affecting the others.

\subsection{Hospital Zone Model}

Hospital devices have very different criticality levels. ICU monitors cant go offline even for a second. Radiology workstations handle large files but can tolerate a brief pause. Front-desk terminals are the least sensitive. We reflected this with three zones:

\begin{itemize}
    \item \textbf{Critical Zone} ($h_1$, $h_2$): Patient monitors, EHR servers. Any downtime is risky.
    \item \textbf{Imaging Zone} ($h_3$, $h_4$): PACS workstations, radiology viewers. Big data transfers but short pauses are fine.
    \item \textbf{Administrative Zone} ($h_5$, $h_6$): Billing, reception, email. Nobody notices a quick hiccup.
\end{itemize}

Each device carries two addresses. One fixed from 10.0.0.0/24---applications bind to this, it never changes. The other from 172.16.0.50--99 is what the rest of the network sees, and it rotates on a timer. The controller rewrites packet headers to translate between these on every frame passing through the switch.

\subsection{SDN Controller}

We implemented the central logic as a Ryu module using OpenFlow~1.3. It handles six functions:

\begin{itemize}
    \item \textbf{IP bootstrapping.} New hosts that broadcast a discovery get assigned a private IP, similar to DHCP.
    \item \textbf{Bidirectional NAT.} Outgoing frames get their source IP swapped to the current rotating address. Incoming frames get their destination restored to the internal one.
    \item \textbf{Layered flow priorities.} ARP at priority 1000, ICMP at 500, per-host NAT at 100, default catch-all at 1. This prevents broad rules from overriding host-specific ones during transitions.
    \item \textbf{Name resolution.} After each rotation, DNS-like hostname mappings get updated so hosts can still reach each other by name.
    \item \textbf{Zone access enforcement.} Before installing any rule, the controller checks whether the source and destination zones are allowed to communicate. Forbidden pairs never get a flow entry.
    \item \textbf{State persistence.} All allocations, rotations, and decisions go into SQLite for audit and crash recovery.
\end{itemize}

An HTTP service on port~8000 exposes endpoints for status, logs, name resolution, and on-demand rotation triggers.

\subsection{Emulated Network}

Six virtual hosts run on Mininet, connected through one Open vSwitch managed by the controller. Each host has a lightweight web server on port~8080 for recieving encrypted messages and sending acknowledgments. A helper daemon on port~8888 allows remote command execution on any host for test automation.

\subsection{Rotation Scheduler}

A standalone Python daemon tracks time per zone and calls the controller's shuffle endpoint over HTTP when any zone's interval expires. Zone membership, timing, and communication permissions all live in a single YAML file (\texttt{policies.yml}).

Access rules follow a simple hierarchy: more critical zones get broader access. Critical hosts can reach anyone. Administrative ones face the strictest limits. Evaluation is top-down, first matching rule wins.

\section{Methodology}

\subsection{Rotation Procedure}

Algorithm~\ref{alg:iphopping} describes the rotation step by step. When a zone's timer expires, the scheduler loops through its hosts, picks an unused IP from the pool, removes old flow entries, installs new ones, updates name mappings, and writes an audit record.

\begin{algorithm}[htbp]
\caption{Zone-Based IP Hopping Procedure}
\label{alg:iphopping}
\begin{algorithmic}[1]
\Require Zone configuration $Z = \{z_1, z_2, z_3\}$, IP pool $P$, interval map $I$
\Ensure Updated NAT mappings and flow rules
\For{each zone $z_i \in Z$}
    \If{$elapsed\_time(z_i) \geq I(z_i)$}
        \For{each host $h_j \in z_i$}
            \State $ip_{old} \gets \text{NAT}[h_j.\text{private\_ip}]$
            \State $ip_{new} \gets \text{RandomSelect}(P \setminus \text{ActiveIPs})$
            \State \text{DeleteFlows}($h_j$, $ip_{old}$)
            \State \text{InstallNATFlows}($h_j$, $ip_{new}$)
            \State \text{NAT}[$h_j.\text{private\_ip}$] $\gets ip_{new}$
            \State \text{UpdateDNS}($h_j.\text{hostname}$, $ip_{new}$)
            \State \text{LogShuffle}($h_j$, $ip_{old}$, $ip_{new}$)
        \EndFor
        \State \text{ResetTimer}($z_i$)
    \EndIf
\EndFor
\end{algorithmic}
\end{algorithm}

Two complementary OpenFlow entries handle the actual rewriting at the switch:

\begin{verbatim}
# Outbound: private_src -> public_src
match = OFPMatch(eth_type=0x0800,
                 ipv4_src=private_ip)
actions = [OFPActionSetField(
               ipv4_src=public_ip),
           OFPActionOutput(out_port)]

# Inbound: public_dst -> private_dst
match = OFPMatch(eth_type=0x0800,
                 ipv4_dst=public_ip)
actions = [OFPActionSetField(
               ipv4_dst=private_ip),
           OFPActionOutput(in_port)]
\end{verbatim}

During transitions, old and new rules coexist briefly so packets already in flight dont get dropped. Old entries are cleaned up only after the switch confirms the new ones are active.

\subsection{Threat Model}

We defined four attacker capabilities:

\begin{itemize}
    \item \textbf{Network scanning.} Discover live hosts and open services using tools like nmap.
    \item \textbf{Lateral movement.} After compromising one host, try reaching higher-value targets in other zones.
    \item \textbf{Address spoofing.} Forge source addresses to exploit trust relationships.
    \item \textbf{Ransomware delivery.} Use network map knowledge to target critical servers.
\end{itemize}

Two hard boundaries apply: the attacker cannot interact with the controller process directly, and the 10.0.0.0/24 internal space is not visible to them. Given these constraints, the defense aims to make recon data expire before its actionable, raise pivot costs, and shrink each device's exposure window.

\subsection{Evaluation Metrics}

We chose four measurable indicators:

\begin{enumerate}
    \item \textbf{Scan yield ($R_A$).} Fraction of live hosts a single scan correctly finds.
    \begin{equation}
        R_A = \frac{|\text{Hosts discovered with valid IPs}|}{|\text{Total active hosts}|}
    \end{equation}

    \item \textbf{Effort multiplier ($E_M$).} Additional scans needed versus a static baseline.
    \begin{equation}
        E_M = \frac{\text{Scans required under MTD}}{\text{Scans required under static}} = \frac{T_{hop}^{-1} \cdot T_{scan}}{1}
    \end{equation}
    $T_{hop}$ = rotation period, $T_{scan}$ = time for one full scan.

    \item \textbf{Delivery rate ($PDR$).} Packets arriving during active rotation.
    \begin{equation}
        PDR = \frac{\text{Packets received}}{\text{Packets sent}} \times 100\%
    \end{equation}

    \item \textbf{Swap latency ($\delta_f$).} Time from shuffle trigger to new rules being active on the switch, in milliseconds.
\end{enumerate}

\section{Implementation}

The prototype consists of:

\begin{itemize}
    \item \textbf{Controller.} Python~3.8, Ryu framework, eventlet~0.30.2, packaged in Docker. OpenFlow~1.3 over TCP.
    \item \textbf{Virtual network.} Mininet~2.3.0 with Open vSwitch~2.13.
    \item \textbf{Host services.} Each host runs a small HTTP server on port~8080 that accepts encrypted payloads, checks integrity, and writes a timestamped receipt.
    \item \textbf{Scheduler.} Standalone script that monitors time and triggers the shuffle API.
    \item \textbf{Dashboard.} Browser-based page (HTML/CSS/JS) polling \texttt{/status} and \texttt{/logs} to display live mappings and rotation history.
    \item \textbf{Attack tools.} Nmap for discovery, Wireshark for packet inspection, cURL for encrypted transfers.
\end{itemize}

Because the controller, scheduler, and HTTP server share mutable state, every write acquires a mutex. All changes persist to SQLite across four tables: \texttt{hosts}, \texttt{nat\_mappings}, \texttt{shuffle\_log}, \texttt{dns\_records}.

For end-to-end validation we used the \texttt{/sim/secure\_transfer} endpoint rather than \texttt{pingall}. The reason is \texttt{pingall} occasionally gives false negatives when NAT rewriting interferes with ICMP replies. The secure transfer path does a full encrypted send-recieve-acknowledge cycle between any two hosts.

\section{Experimental Setup}

\subsection{Lab Environment}

Six Mininet hosts ($h_1$--$h_6$) connect through one Open vSwitch under a single Ryu controller. $h_1, h_2$ are Critical, $h_3, h_4$ are Imaging, $h_5, h_6$ are Administrative. Fixed IPs range from 10.0.0.1 to 10.0.0.6. The rotation pool covers 172.16.0.50--99 (50 addresses for 6 hosts). The controller container uses \texttt{--network host} to share the namespace with Mininet.

\subsection{Rotation Timing}

Table~\ref{tab:zoneconfig} lists the intervals after iterative tuning.

\begin{table}[htbp]
\centering
\caption{Zone-Based IP Hopping Configuration}
\label{tab:zoneconfig}
\begin{tabular}{lccc}
\toprule
\textbf{Hospital Zone} & \textbf{Risk Level} & \textbf{Hosts} & \textbf{Interval (s)} \\
\midrule
Administrative & Low & $h_5, h_6$ & 80 \\
Imaging & Medium & $h_3, h_4$ & 100 \\
Critical & High & $h_1, h_2$ & 120 \\
\bottomrule
\end{tabular}
\end{table}

It seems counterintuitive that Critical rotates slowest. Early experiments explained why: when Critical hosts rotated too quickly, a timing gap appeared between the controller installing new rules and the host daemon updating its bindings. That gap broke active telemetry streams. Since these devices may be keeping patients alive, we opted for a longer but perfectly reliable cycle. Admin machines tolerate momentary blips, so they get the fastest rate.

\subsection{Test Scenarios}

Three experiments were conducted:

\begin{enumerate}
    \item \textbf{Scan degradation.} Repeated \texttt{nmap -sn 172.16.0.0/24} sweeps to track how the attacker's view decays with each rotation.
    \item \textbf{Pivot attempts.} From an Administrative host, identify a Critical target, wait, then try to connect. The IP may have rotated in between.
    \item \textbf{Legitimate traffic under rotation.} Host pairs exchanged encrypted payloads through the REST endpoint during full rotation cycles. Successes and failures were counted.
\end{enumerate}

\textit{Assumptions.} (i)~Attacker has no access to the controller or internal IPs. (ii)~One scan takes roughly one rotation period. (iii)~Mininet behavior approximates physical OpenFlow switches. (iv)~Single-switch latency values are a lower bound.

\section{Results and Discussion}

\subsection{Scan Accuracy Drops Sharply}

We ran 20 consecutive scans (15--20\,s each) against both a static and a rotating network.

On the static network every scan returned all six hosts, as expected. With MTD active:

\begin{itemize}
    \item \textbf{Scan 1}: All six visible---rotation hadnt started yet.
    \item \textbf{Scans 2--5}: Average of 3.2 hosts per scan as Administrative IPs began cycling.
    \item \textbf{Scans 6--20}: Roughly one host per scan once every zone had rotated at least once.
\end{itemize}

Steady-state yield settled at \textbf{16.7\%} (one out of six). Thats an \textbf{83.3\% reduction} from the static baseline. The single visible host is whichever was assigned its address most recently.

This means scan data has a very short shelf life. An attacker who scans once gets a map that expires within 80--120 seconds depending on the zone. Repeated scanning gives diminishing returns because each new sweep mostly captures addresses that will rotate again soon. From a defender perspective, this is exactly what we wanted: intelligence that expires faster than the attacker can act on it.

\subsection{Attacker Effort Goes Up Substantially}

On a static network, one scan reveals everything. With rotation, each scan expires at the next cycle. For Administrative (80\,s interval, 18\,s scan):

\begin{equation}
    E_M = \frac{80}{18} \approx 4.4\times \text{ (Administrative zone)}
\end{equation}

For Critical at 120\,s:

\begin{equation}
    E_M = \frac{120}{18} \approx 6.7\times \text{ (Critical zone)}
\end{equation}

The weighted average is \textbf{4.7$\times$}. So the attacker needs almost five times more scanning to maintain the same level of awareness.

Worth noting: this is a conservative estimate. In practice the attacker also has to correlate results across scans to figure out which addresses are still valid, adding more overhead on top of just the raw scan count.

\subsection{Lateral Movement Mostly Blocked}

We ran 15 pivot experiments. Each trial: scan the target, pause briefly, attempt connection.

\begin{itemize}
    \item \textbf{Static}: 15/15 succeeded (100\%).
    \item \textbf{MTD}: 3/15 succeeded (\textbf{20\%}).
\end{itemize}

12 failures happened because the target address had rotated before the connection attempt. The zone access policy added an independent barrier---Admin-to-Critical traffic gets blocked regardless of whether the address is right.

This creates a dual-defense effect. Even when rotation alone might not be fast enough (the 3 successes), zone policies catch what slips through. Both layers have to fail simultaneously for an attack to land, which is a much harder condition for the attacker.

\subsection{Operational Overhead Remains Low}

Table~\ref{tab:perfeval} summarizes all measured metrics.

\begin{table}[htbp]
\centering
\caption{Performance Evaluation: Static vs.\ MTD-HealthNet}
\label{tab:perfeval}
\begin{tabular}{lccc}
\toprule
\textbf{Metric} & \textbf{Static} & \textbf{MTD} & \textbf{Impact} \\
\midrule
Packet Delivery Ratio & 99.8\% & 97.8\% & \textbf{2.0\% decrease} \\
Avg.\ RTT (steady) & 0.34 ms & 0.41 ms & +0.07 ms \\
Avg.\ RTT (during hop) & --- & 2.8 ms & Transient spike \\
Flow Update Latency & --- & 12.4 ms & Acceptable \\
Recon.\ Accuracy & 100\% & 16.7\% & \textbf{83.3\% reduction} \\
Lateral Movement Success & 100\% & 20\% & \textbf{80\% reduction} \\
Attack Effort Multiplier & 1$\times$ & 4.7$\times$ & \textbf{4.7$\times$ increase} \\
\bottomrule
\end{tabular}
\end{table}

Breaking down the operational cost:

\begin{itemize}
    \item The \textbf{2\% delivery drop} only occurred during the narrow moment when old rules were being replaced. The overlap window kept it from getting worse.
    \item Baseline RTT increased by \textbf{0.07~ms} from per-packet header rewriting. During active swaps it spiked to 2.8~ms but settled back within a second or two.
    \item \textbf{12.4~ms swap latency} is invisible to clinical applications that work on timescales of whole seconds.
    \item Encrypted payload exchanges through the REST endpoint succeeded \textbf{100\%} of the time, including during rotation events.
\end{itemize}

Giving up 2\% delivery and 0.07~ms latency to get 83.3\% less scan accuracy and 80\% fewer successful pivots seems like a good deal. Hospital informatics cares about seconds, not sub-millisecond increments. And since rotation timing is per-zone, administrators keep fine-grained control---a neonatal ICU could use 180-second cycles while guest wifi runs at 45 seconds within the same framework.

\section{Case Studies}

\subsection{Scenario A: Ransomware Targeting Patient Records}

\textit{Setup.} Attacker compromises $h_5$ (Administrative) and targets $h_1$ (Critical, storing patient records).

\textit{Static network.} One scan reveals $h_1$'s address. The attacker crafts the payload at leisure and delivers it.

\textit{With MTD-HealthNet.} Scan shows $h_1$ at 172.16.0.57. Weaponization takes 60--90 seconds. But $h_1$ cycles every 120\,s, so 172.16.0.57 probably belongs to someone else or nobody by the time the payload ships. Even rescanning wont help because the access policy independently blocks Admin-to-Critical traffic.

\subsection{Scenario B: Intra-Zone Movement in Imaging}

\textit{Setup.} Attacker on $h_3$ (Imaging) wants to exfiltrate radiology data from $h_4$.

\textit{With MTD-HealthNet.} $h_4$ is at 172.16.0.62. One Imaging cycle later (100\,s), its at 172.16.0.78. The old address leads nowhere. A fresh scan might find the new one, but another rotation is already approaching. The attacker falls into a loop where intelligence expires almost as fast as it can be gathered.

\subsection{Scenario C: Continuous Reconnaissance}

\textit{Setup.} External attacker runs nonstop scan sweeps.

\textit{With MTD-HealthNet.} First sweep gets six addresses. After 80\,s, Administrative entries have moved. At 100\,s, Imaging follows. By 120\,s everything has cycled at least once. Accumulated data is a mix of valid and stale IPs---maybe one in six is still current. Moreover, every observed address comes from the 172.16.0.0/16 pool, revealing nothing about the real 10.0.0.0/24 layout.

\subsection{Limitations}

\begin{itemize}
    \item \textbf{Single point of failure.} One Ryu process handles everything. If it crashes, all defensive functions stop. No redundancy exists yet.
    \item \textbf{Simplified topology.} A single virtual switch doesnt represent real multi-building hospital networks accurately.
    \item \textbf{Legacy device gaps.} Equipment with hardcoded peer addresses would need a translation proxy we havent implemented.
    \item \textbf{Limited entropy.} 50 candidate IPs for 6 devices yields about $\log_2\binom{50}{6} \approx 22.7$ bits per cycle. Useful but far below cryptographic standards.
    \item \textbf{Minimum rotation speed.} Below roughly 30 seconds, timing mismatches between rule installation and daemon updates degraded reliability.
\end{itemize}

\section{Conclusion}

We set out to test whether regular IP rotation is practical for hospital networks. MTD-HealthNet reduced scan effectiveness by 83.3\%, increased attacker scanning workload by 4.7$\times$, and blocked 80\% of cross-zone pivot attempts. Delivery rates stayed above 97.8\% with negligible added latency.

No special hardware is required. Zone-level timing lets security teams balance protection and availability per department.

\section{Future Work}

Several directions remain open:

\begin{itemize}
    \item \textbf{Adaptive timing.} Replace fixed intervals with anomaly-driven cadences---speed up under suspected attack, slow down during quiet periods.
    \item \textbf{IDS-triggered rotation.} Let an intrusion detection module force an immediate IP swap for flagged hosts, adding reactive defense on top of scheduled rotation.
    \item \textbf{Multi-switch scaling.} Expand the testbed to multiple interconnected switches and subnets, profiling controller resource use as device count grows.
    \item \textbf{Port mutation.} Rotate listening ports alongside addresses, rewriting them transparently at the switch so endpoint software stays unaware.
    \item \textbf{Game-theoretic optimization.} Model the defender-attacker interaction formally to derive mathematically optimal intervals for given scan speeds and network sizes.
    \item \textbf{Hospital pilot.} Deploy on physical hospital infrastructure to collect measurements under production traffic and regulatory constraints.
\end{itemize}

\begin{thebibliography}{00}

\bibitem{dishington2019ipmux}
C.~Dishington, D.~P.~Sharma, D.~S.~Kim, J.-H.~Cho, T.~J.~Moore, and F.~F.~Nelson, ``Security and performance assessment of IP multiplexing moving target defence in software defined networks,'' in \textit{Proc.\ 2019 18th IEEE Int.\ Conf.\ Trust, Security and Privacy in Computing and Communications / 13th IEEE Int.\ Conf.\ Big Data Science and Engineering (TrustCom/BigDataSE)}, IEEE, 2019, doi: 10.1109/trustcom/bigdatase.2019.00046.

\bibitem{zal2024lossless}
M.~\.{Z}al, M.~Michalski, and P.~Zwierzykowski, ``Implementation of a lossless moving target defense mechanism,'' \textit{Electronics}, MDPI AG, 2024, doi: 10.3390/electronics13050918.

\bibitem{lei2018survey}
C.~Lei, H.-Q.~Zhang, J.-L.~Tan, Y.-C.~Zhang, and X.-H.~Liu, ``Moving target defense techniques: A survey,'' \textit{Security and Communication Networks}, Wiley, 2018, doi: 10.1155/2018/3759626.

\bibitem{cai2016mtd}
G.-L.~Cai, B.-S.~Wang, W.~Hu, and T.-Z.~Wang, ``Moving target defense: State of the art and characteristics,'' \textit{Frontiers of Information Technology \& Electronic Engineering}, Zhejiang University Press, 2016, doi: 10.1631/fitee.1601321.

\bibitem{huang2021legacy}
C.-W.~Huang, I.-H.~Liu, J.-S.~Li, C.-C.~Wu, C.-F.~Li, and C.-G.~Liu, ``A legacy infrastructure-based mechanism for moving target defense,'' in \textit{Proc.\ 2021 IEEE 3rd Eurasia Conf.\ Biomedical Engineering, Healthcare and Sustainability (ECBIOS)}, IEEE, 2021, doi: 10.1109/ecbios51820.2021.9510261.

\end{thebibliography}

\end{document}
